\begin{abstract}
	This study investigates the applicability of the Beardwood-Halton-Hammersley (BHH) formula for estimating Traveling Salesman Problem (TSP) path lengths in real-world delivery scenarios using OpenStreetMap data from various Dutch neighborhoods.
	While the BHH formula yields accurate predictions with a Mean Absolute Percentage Error (MAPE) of 6.6\% when the area-specific proportionality constant $\beta$ is known, its performance degrades significantly (with MAPE exceeding 22\%) when using a nationwide average $\beta$.
	This variability limits its practical utility, since this formula is mainly used in applications where the true value of $\beta$ is not known.
	Even with area-specific estimates of $\beta$, accuracy varies depending on address distribution and road network connectivity.
	To overcome these limitations, a supervised learning model leveraging geographic and spatial features was developed, achieving improved generalizability and a lower MAPE of 10.6\% on the test set.
	However, prediction errors remain high in areas with clustered or irregular address distribution.
\end{abstract}
